% ============================================================
%  PROBABILITY PROBLEMS — UNIT 1
%  One file per unit; add \day sections as the course advances.
%  Compile with pdflatex (works on Overleaf out of the box).
% ============================================================

\documentclass[12pt, a4paper]{article}

% ---- encoding & language -----------------------------------
\usepackage[utf8]{inputenc}
\usepackage[T1]{fontenc}
\usepackage[english]{babel}        % change to [spanish] if needed

% ---- mathematics -------------------------------------------
\usepackage{amsmath, amssymb, amsthm}

% ---- R code listings ---------------------------------------
\usepackage{listings}
\usepackage{xcolor}

\definecolor{codegreen}{rgb}{0.2,0.5,0.2}
\definecolor{codegray}{rgb}{0.5,0.5,0.5}
\definecolor{codeblue}{rgb}{0.1,0.1,0.8}
\definecolor{backcolor}{rgb}{0.97,0.97,0.97}

\lstdefinestyle{Rstyle}{
  language=R,
  backgroundcolor=\color{backcolor},
  basicstyle=\ttfamily\small,
  keywordstyle=\color{codeblue}\bfseries,
  commentstyle=\color{codegreen}\itshape,
  stringstyle=\color{orange},
  numberstyle=\tiny\color{codegray},
  numbers=left,
  stepnumber=1,
  frame=single,
  rulecolor=\color{codegray},
  breaklines=true,
  showstringspaces=false,
  tabsize=2,
  captionpos=b
}
\lstset{style=Rstyle}

% ---- layout ------------------------------------------------
\usepackage{geometry}
\geometry{margin=2.5cm}

\usepackage{enumitem}             % better lists
\usepackage{parskip}              % space between paragraphs, no indent
\usepackage{titlesec}             % section formatting
\usepackage{fancyhdr}             % header / footer
\usepackage{lastpage}             % "Page X of Y"
\usepackage{hyperref}             % clickable links (keep last)

% ---- header & footer ---------------------------------------
\pagestyle{fancy}
\fancyhf{}
\lhead{\textbf{Probability} — Unit 1: \unitname}
\rhead{\today}
\cfoot{Page \thepage\ of \pageref{LastPage}}
\renewcommand{\headrulewidth}{0.4pt}

% ---- problem counter ---------------------------------------
%  Resets automatically with each \day command below.
\newcounter{problem}[section]
\renewcommand{\theproblem}{\thesection.\arabic{problem}}

\newcommand{\problem}[1][]{%          \problem  or  \problem[title]
  \refstepcounter{problem}%
  \medskip
  \noindent\textbf{Problem~\theproblem%
    \ifx&#1&\else~—~#1\fi}%          optional title after the dash
  \par\smallskip
}

% ---- \day shortcut -----------------------------------------
%  Usage:  \day{1}{Descriptive subtitle}
\newcommand{\day}[2]{%
  \section*{Day #1 \normalfont\small— #2}%
  \addcontentsline{toc}{section}{Day #1: #2}%
  % Reset problem counter relative to day, not section number:
  \setcounter{problem}{0}%
}

% ---- unit name (edit here each time) -----------------------
\newcommand{\unitname}{Sample Spaces and Events}

% ============================================================
\begin{document}
% ============================================================

\begin{center}
  {\LARGE \textbf{Unit 1: \unitname}} \\[0.4em]
  {\large Probability — Problem Sets} \\[0.2em]
  {\small Last updated: \today}
\end{center}

\tableofcontents
\bigskip\hrule\bigskip

% ============================================================
\day{1}{Introduction to Sample Spaces}
% ============================================================

\problem
A fair six-sided die is rolled once.
\begin{enumerate}[label=(\alph*)]
  \item Write out the sample space $\Omega$.
  \item Let $A$ be the event ``the result is even''. List the elements of $A$.
  \item Compute $P(A)$.
\end{enumerate}

\problem[Complement rule]
Let $B$ be the event ``the result is greater than 4''.
\begin{enumerate}[label=(\alph*)]
  \item Find $P(B)$ and $P(B^c)$.
  \item Verify that $P(B) + P(B^c) = 1$.
\end{enumerate}

\problem[Simulation in R]
The following code simulates 10{,}000 rolls of a fair die and estimates $P(A)$
from Problem~1.1.
\begin{lstlisting}
set.seed(42)
rolls <- sample(1:6, size = 10000, replace = TRUE)

# Estimate P(even)
mean(rolls %% 2 == 0)
\end{lstlisting}
Run the code, compare with the theoretical value, and explain any discrepancy.

% ============================================================
\day{2}{Operations on Events}
% ============================================================

\problem
Let $\Omega = \{1, 2, 3, 4, 5, 6\}$, $A = \{2, 4, 6\}$, $B = \{4, 5, 6\}$.
Compute:
\begin{enumerate}[label=(\alph*)]
  \item $A \cup B$
  \item $A \cap B$
  \item $A \setminus B$
  \item $A^c \cap B$
\end{enumerate}

\problem
Using the same events, verify the inclusion–exclusion formula
\[
  P(A \cup B) = P(A) + P(B) - P(A \cap B).
\]

\problem[De Morgan's laws]
Prove or illustrate with an example that
\[
  (A \cup B)^c = A^c \cap B^c.
\]

% ============================================================
\day{3}{Conditional Probability}
% ============================================================

\problem
A card is drawn at random from a standard 52-card deck.
Let $A = \{\text{the card is a heart}\}$ and $B = \{\text{the card is a face card}\}$.
\begin{enumerate}[label=(\alph*)]
  \item Compute $P(A \mid B)$.
  \item Are $A$ and $B$ independent? Justify.
\end{enumerate}

\problem[Simulation in R]
Estimate $P(A \mid B)$ by simulation and compare with your answer above.
\begin{lstlisting}
set.seed(7)
suits  <- rep(c("hearts","diamonds","clubs","spades"), each = 13)
values <- rep(c(2:10, "J","Q","K","A"), times = 4)
deck   <- data.frame(suit = suits, value = values)

n <- 100000
draws <- deck[sample(nrow(deck), n, replace = TRUE), ]

face_cards <- draws$value %in% c("J", "Q", "K")
hearts     <- draws$suit == "hearts"

# Estimate P(heart | face card)
mean(hearts[face_cards])
\end{lstlisting}

% ============================================================
%  Add more \day{4}{...} blocks here as the unit progresses.
% ============================================================

\end{document}